\chapter{Reações pericíclicas}\label{ch:b3:pericyclic}

A luz do sol que incide sobre a pele abre um anel do 7-desidrocolesterol em uma fração
de segundo; o calor do corpo desloca então um átomo de hidrogênio através de sete átomos de carbono,
e o produto é a vitamina D$_3$. Nenhuma das duas etapas tem intermediário: ligações se rompem
e se formam juntas, ao redor de um anel de átomos, ao longo de um único estado de transição. A
primeira etapa só ocorre com luz e a segunda só com calor, e cada uma dá
um único estereoisômero. Este capítulo explica por quê, a partir da simetria dos
orbitais moleculares envolvidos: as regras encontradas por Robert Woodward e Roald
Hoffmann, que preveem se uma reação desse tipo pode ocorrer e o que ela dá.

\begin{recall}
O volume do 2.º ano deu os coeficientes de Hückel dos polienos, os anéis aromáticos e
antiaromáticos, a HOMO, a LUMO e os orbitais de fronteira, as reações concertadas e
a reação de Diels--Alder de um dieno com um dienófilo, com sua regra endo; o
volume do 1.º ano, os descritores \textit{Z}/\textit{E}. O \cref{ch:b3:point-groups} definiu o eixo $C_2$
e o \omterm{def:b3:point-groups:mirror}{plano especular}, o \cref{ch:b3:photochemistry}, os estados excitados e a fotocicloadição
[2+2].
\end{recall}

\begin{omfigure}
\includegraphics[width=0.55\linewidth]{images/book4/ai/b3-26-terrace.jpg}

{\small Lendo em um terraço ensolarado: a luz ultravioleta que chega à pele inicia
a síntese da vitamina D$_3$ por uma abertura de anel pericíclica.}
\end{omfigure}

\section{Quatro famílias}

\begin{definition}[Reação pericíclica]\label{def:b3:pericyclic:pericyclic}
Uma \emph{reação pericíclica}\index{reacao periciclica@reação pericíclica} é uma reação concertada
cujo estado de transição é um arranjo cíclico de orbitais em interação, no qual todas
as ligações se formam e se rompem ao redor do anel ao mesmo tempo.
\end{definition}

\begin{definition}[Famílias de reações pericíclicas]\label{def:b3:pericyclic:families}
Em uma \emph{reação eletrocíclica}\index{reacao eletrociclica@reação eletrocíclica}, um polieno
conjugado se fecha em um anel formando uma ligação $\sigma$ entre suas extremidades, ou o
inverso. Em uma \emph{cicloadição}\index{cicloadicao@cicloadição}, dois sistemas $\pi$ se unem
em um anel formando duas ligações $\sigma$; uma cicloadição $[m + n]$ une componentes
de $m$ e de $n$ átomos (ou elétrons). Em um \emph{rearranjo
sigmatrópico}\index{rearranjo sigmatropico@rearranjo sigmatrópico}, uma ligação $\sigma$ migra ao longo de um sistema $\pi$;
em um deslocamento $[i, j]$, a ligação passa dos átomos 1 e 1 para os átomos $i$ e $j$
dos dois fragmentos. Uma \emph{reação queletrópica}\index{reacao queletropica@reação queletrópica}
forma ou rompe duas ligações com o mesmo átomo, como quando o dióxido de enxofre
deixa um sulfoleno. Em uma \emph{reação eno}\index{reacao eno@reação eno}, um alqueno
com um hidrogênio alílico se adiciona a uma ligação $\pi$, e o hidrogênio migra com ele.
\end{definition}

\begin{definition}[Suprafacial e antarafacial]\label{def:b3:pericyclic:faciality}
Um componente reage de modo \emph{suprafacial}\index{suprafacial} quando suas duas novas
ligações se formam, ou suas duas antigas se rompem, na mesma face de seu sistema $\pi$ (ou com
retenção em uma ligação $\sigma$), e de modo \emph{antarafacial}\index{antarafacial} quando elas
ficam em faces opostas (ou com inversão).
\end{definition}

\section{Reações eletrocíclicas}

\begin{definition}[Conrotatório e disrotatório]\label{def:b3:pericyclic:rotation-modes}
Em um fechamento eletrocíclico de anel, os dois grupos terminais giram em torno das ligações
que os unem à cadeia. O movimento é \emph{conrotatório}\index{conrotatorio@conrotatório}
quando ambos giram no mesmo sentido (ambos no sentido horário, vistos ao longo da cadeia), e
\emph{disrotatório}\index{disrotatorio@disrotatório} quando giram em sentidos opostos.
\end{definition}

\begin{theorem}[Reações eletrocíclicas]\label{thm:b3:pericyclic:electrocyclic}
Uma \omterm{def:b3:pericyclic:families}{reação eletrocíclica} térmica de um polieno com $4n$ elétrons $\pi$ é
\omterm{def:b3:pericyclic:rotation-modes}{conrotatória}, e com $4n + 2$ elétrons $\pi$, \omterm{def:b3:pericyclic:rotation-modes}{disrotatória}.
\end{theorem}

\begin{proof}
A nova ligação $\sigma$ se forma a partir dos orbitais p terminais da HOMO, que na
reação térmica contém os elétrons de maior energia; ela só é ligante se os
lobos que se voltam um para o outro têm o mesmo sinal. Pela fórmula de Hückel
(\cref{ch:b3:solid-state}), os coeficientes do orbital $k$ de uma cadeia de $N$ átomos são
proporcionais a $\sin(jk\pi/(N + 1))$, de modo que as duas extremidades, $j = 1$ e $j = N$, têm
coeficientes $\sin(k\pi/(N + 1))$ e $\sin(Nk\pi/(N + 1)) = (-1)^{k+1}\sin(k\pi/(N + 1))$: o mesmo sinal para $k$ ímpar,
sinais opostos para $k$ par. Com $N$ átomos e $N$ elétrons, a HOMO é
$k = N/2$. Para $N = 4n$, $k = 2n$ é par: os lobos superiores das extremidades têm sinais
opostos, e as duas extremidades precisam girar no mesmo sentido, cada uma trazendo para dentro seu lobo
do sinal correspondente: \omterm{def:b3:pericyclic:rotation-modes}{conrotatório}. Para $N = 4n + 2$, $k = 2n + 1$ é ímpar: os lobos superiores
têm o mesmo sinal, e voltá-los um para o outro é \omterm{def:b3:pericyclic:rotation-modes}{disrotatório}.
\end{proof}

\begin{omfigure}
\begin{tikzpicture}[font=\small, >=Stealth]
  % butadiene HOMO psi2: 0.60 0.37 -0.37 -0.60 (lobes scaled x1.6)
  \foreach \x/\c in {0/0.60, 1/0.37, 2/-0.37, 3/-0.60} {
    \pic at (\x,0) {omp={1.6*\c}{90}};
    \fill (\x,0) circle (1.2pt);
  }
  \draw[black!60] (0,0) -- (3,0);
  \draw[->, thick, omThm] (-0.45,0.55) arc[start angle=130, end angle=20, radius=0.7];
  \draw[->, thick, omThm] (2.55,0.55) arc[start angle=130, end angle=20, radius=0.7];
  \node[font=\scriptsize, align=center] at (1.5,-1.3) {butadieno, $\psi_2$ (HOMO, 4 elétrons):\\ extremidades de sinais opostos: \textbf{conrotatório}};
  \begin{scope}[xshift=6.2cm]
  % hexatriene HOMO psi3: 0.52 0.23 -0.42 -0.42 0.23 0.52
  \foreach \x/\c in {0/0.52, 0.8/0.23, 1.6/-0.42, 2.4/-0.42, 3.2/0.23, 4.0/0.52} {
    \pic at (\x,0) {omp={1.6*\c}{90}};
    \fill (\x,0) circle (1.2pt);
  }
  \draw[black!60] (0,0) -- (4.0,0);
  \draw[->, thick, omThm] (-0.45,0.55) arc[start angle=130, end angle=20, radius=0.7];
  \draw[->, thick, omThm] (4.45,0.55) arc[start angle=50, end angle=160, radius=0.7];
  \node[font=\scriptsize, align=center] at (2.0,-1.3) {hexatrieno, $\psi_3$ (HOMO, 6 elétrons):\\ extremidades de mesmo sinal: \textbf{disrotatório}};
  \end{scope}
\end{tikzpicture}

{\small A HOMO do butadieno e a do hexatrieno, com os lobos em escala segundo seus
coeficientes de Hückel (azul e laranja: os dois sinais), desenhadas ao longo da cadeia.
Para que uma ligação se forme entre as extremidades, os lobos que se voltam um para o outro
precisam ter o mesmo sinal: as duas extremidades giram no sentido horário no butadieno, em
sentidos opostos no hexatrieno (setas vermelhas).}
\end{omfigure}

A regra é estereoespecífica. O (\textit{E},\textit{E})-hexa-2,4-dieno, aquecido, se fecharia
de modo \omterm{def:b3:pericyclic:rotation-modes}{conrotatório} no \textit{trans}-3,4-dimetilciclobuteno; na prática, é o
ciclobuteno tensionado que se abre, e o \textit{cis}-3,4-dimetilciclobuteno se abre
de modo \omterm{def:b3:pericyclic:rotation-modes}{conrotatório} apenas no (\textit{E},\textit{Z})-hexa-2,4-dieno. O (\textit{E},\textit{Z},\textit{E})-octa-2,4,6-trieno se fecha
de modo \omterm{def:b3:pericyclic:rotation-modes}{disrotatório} no \textit{cis}-5,6-dimetilciclo-hexa-1,3-dieno.

\begin{proposition}[Reações eletrocíclicas fotoquímicas]\label{prop:b3:pericyclic:photo-electrocyclic}
Sob luz, as regras se invertem: $4n$ elétrons, \omterm{def:b3:pericyclic:rotation-modes}{disrotatório}; $4n + 2$, \omterm{def:b3:pericyclic:rotation-modes}{conrotatório}.
\end{proposition}

\begin{proof}
A absorção de um fóton promove um elétron da HOMO $k = N/2$ para a LUMO $k = N/2 + 1$
(\cref{ch:b3:photochemistry}), que passa a ser o orbital ocupado mais alto e
controla as extremidades. Seu $k$ tem a paridade oposta, de modo que o sinal relativo de seus
coeficientes terminais se inverte, e com ele o sentido de rotação.
\end{proof}

\begin{method}[Prever a estereoquímica de uma reação eletrocíclica]\label{met:b3:pericyclic:predict-stereo}
\begin{enumerate}
  \item Conte os elétrons $\pi$ do polieno de cadeia aberta ($4n$ ou $4n + 2$); anote calor ou
        luz.
  \item Leia o modo: térmico, $4n$ con, $4n + 2$ dis; fotoquímico, o inverso.
  \item Desenhe as extremidades com seus substituintes, na conformação que
        se fecha (s-cis); gire-as de 90° no modo escolhido.
  \item Veja se os dois substituintes terminam na mesma face do anel
        (\textit{cis}) ou em faces opostas (\textit{trans}). Para uma abertura de anel, faça a mesma
        análise em sentido inverso.
\end{enumerate}
\end{method}

\section{Diagramas de correlação de orbitais}

O argumento original de Woodward e Hoffmann considera todos os orbitais, não apenas
a HOMO. Ao longo de um caminho \omterm{def:b3:pericyclic:rotation-modes}{conrotatório}, a molécula conserva um eixo $C_2$; ao longo de um
caminho \omterm{def:b3:pericyclic:rotation-modes}{disrotatório}, um \omterm{def:b3:point-groups:mirror}{plano especular} $\sigma$. Cada orbital do reagente e do
produto é simétrico (S) ou antissimétrico (A) em relação ao elemento que
se conserva.

\begin{definition}[Diagrama de correlação de orbitais]\label{def:b3:pericyclic:orbital-correlation}
Um \emph{diagrama de correlação de orbitais}\index{diagrama de correlacao de orbitais@diagrama de correlação de orbitais} liga
cada orbital do reagente ao orbital do produto de mesma simetria,
em relação aos \omterm{def:b3:point-groups:operation}{elementos de simetria} conservados ao longo do caminho de reação, tomando-os
em ordem de energia: o S mais baixo com o S mais baixo, e assim por diante.
\end{definition}

\begin{proposition}[Conservação da simetria orbital]\label{prop:b3:pericyclic:correlation}
Se todo orbital ocupado no reagente se correlaciona com um orbital ocupado no
produto, a reação térmica é permitida; se um orbital ocupado se correlaciona
com um vazio, ela é proibida, com uma barreira alta.
\end{proposition}

\begin{proof}[Argumento]
Ao longo do caminho, um orbital conserva sua simetria e sua energia varia
continuamente; orbitais de mesma simetria não podem se cruzar (eles se misturam e se repelem),
os de simetrias diferentes podem. Quando um orbital duplamente ocupado do reagente
precisa se tornar um orbital antiligante do produto, o estado fundamental do
reagente leva a um estado duplamente excitado do produto, e a energia do estado fundamental
sobe abruptamente até que, já no fim do caminho, os estados de mesma simetria
se misturem: a barreira é alta.
\end{proof}

\begin{omfigure}
\begin{tikzpicture}[font=\scriptsize]
  \foreach \px/\title/\el in {0/disrotatório: $\sigma$ conservado/s, 7.6/conrotatório: $C_2$ conservado/c} {
    \begin{scope}[xshift=\px cm]
    \node[font=\small\bfseries] at (2.0,5.1) {\title};
    \node at (0,4.7) {butadieno};
    \node at (4.0,4.7) {ciclobuteno};
    \end{scope}
  }
  % ---- disrotatory (mirror): butadiene psi1 S, psi2 A, psi3 S, psi4 A;
  %      cyclobutene sigma S, pi S, pi* A, sigma* A
  \pic (b1) at (0,0.6) {omlevel={ud}{1.0}}; \node[anchor=east] at (-0.6,0.6) {$\psi_1$ S};
  \pic (b2) at (0,1.8) {omlevel={ud}{1.0}}; \node[anchor=east] at (-0.6,1.8) {$\psi_2$ A};
  \pic (b3) at (0,3.0) {omlevel={}{1.0}}; \node[anchor=east] at (-0.6,3.0) {$\psi_3$ S};
  \pic (b4) at (0,4.1) {omlevel={}{1.0}}; \node[anchor=east] at (-0.6,4.1) {$\psi_4$ A};
  \pic (c1) at (4.0,0.1) {omlevel={ud}{1.0}}; \node[anchor=west] at (4.6,0.1) {$\sigma$ S};
  \pic (c2) at (4.0,1.3) {omlevel={ud}{1.0}}; \node[anchor=west] at (4.6,1.3) {$\pi$ S};
  \pic (c3) at (4.0,3.4) {omlevel={}{1.0}}; \node[anchor=west] at (4.6,3.4) {$\pi^*$ A};
  \pic (c4) at (4.0,4.4) {omlevel={}{1.0}}; \node[anchor=west] at (4.6,4.4) {$\sigma^*$ A};
  \draw[omcorr] (b1-r) -- (c1-l);
  \draw[omcorrforbidden] (b2-r) -- (c3-l);
  \draw[omcorr] (b3-r) -- (c2-l);
  \draw[omcorr] (b4-r) -- (c4-l);
  % ---- conrotatory (C2): butadiene psi1 A, psi2 S, psi3 A, psi4 S;
  %      cyclobutene sigma S, pi A, pi* S, sigma* A
  \begin{scope}[xshift=7.6cm]
  \pic (d1) at (0,0.6) {omlevel={ud}{1.0}}; \node[anchor=east] at (-0.6,0.6) {$\psi_1$ A};
  \pic (d2) at (0,1.8) {omlevel={ud}{1.0}}; \node[anchor=east] at (-0.6,1.8) {$\psi_2$ S};
  \pic (d3) at (0,3.0) {omlevel={}{1.0}}; \node[anchor=east] at (-0.6,3.0) {$\psi_3$ A};
  \pic (d4) at (0,4.1) {omlevel={}{1.0}}; \node[anchor=east] at (-0.6,4.1) {$\psi_4$ S};
  \pic (e1) at (4.0,0.1) {omlevel={ud}{1.0}}; \node[anchor=west] at (4.6,0.1) {$\sigma$ S};
  \pic (e2) at (4.0,1.3) {omlevel={ud}{1.0}}; \node[anchor=west] at (4.6,1.3) {$\pi$ A};
  \pic (e3) at (4.0,3.4) {omlevel={}{1.0}}; \node[anchor=west] at (4.6,3.4) {$\pi^*$ S};
  \pic (e4) at (4.0,4.4) {omlevel={}{1.0}}; \node[anchor=west] at (4.6,4.4) {$\sigma^*$ A};
  \draw[omcorr] (d1-r) -- (e2-l);
  \draw[omcorr] (d2-r) -- (e1-l);
  \draw[omcorr] (d3-r) -- (e4-l);
  \draw[omcorr] (d4-r) -- (e3-l);
  \end{scope}
\end{tikzpicture}

{\small \omterm{def:b3:pericyclic:orbital-correlation}{Diagramas de correlação de orbitais} para butadieno $\rightleftharpoons$ ciclobuteno. À esquerda, o
caminho \omterm{def:b3:pericyclic:rotation-modes}{disrotatório} conserva um \omterm{def:b3:point-groups:mirror}{plano especular}: a $\psi_2$ ocupada (A) se correlaciona com o
$\pi^*$ vazio (A) do ciclobuteno (tracejado vermelho): termicamente proibido. À direita, o
caminho \omterm{def:b3:pericyclic:rotation-modes}{conrotatório} conserva um eixo $C_2$: os orbitais ocupados vão para orbitais ocupados:
termicamente permitido.}
\end{omfigure}

\begin{method}[Construir um diagrama de correlação]\label{met:b3:pericyclic:correlation-diagram}
\begin{enumerate}
  \item Encontre os \omterm{def:b3:point-groups:operation}{elementos de simetria} que cortam ao meio as ligações formadas e rompidas e que se
        conservam ao longo de todo o caminho escolhido.
  \item Liste os orbitais participantes, do reagente e do produto, em ordem de
        energia, e rotule cada um como S ou A para cada elemento.
  \item Ligue o orbital mais baixo de cada tipo de simetria de um lado ao mais baixo
        do mesmo tipo do outro, depois o seguinte, sem nunca cruzar duas linhas de
        mesma simetria.
  \item Verifique os orbitais ocupados: todos para ocupados, permitido; um para vazio,
        proibido.
\end{enumerate}
\end{method}

\begin{definition}[Diagrama de correlação de estados]\label{def:b3:pericyclic:state-correlation}
Um \emph{diagrama de correlação de estados}\index{diagrama de correlacao de estados@diagrama de correlação de estados} liga os
estados eletrônicos do reagente e do produto, cada um construído a partir das configurações
orbitais e rotulado por sua simetria global, ao longo do caminho; estados de
mesma simetria não se cruzam.
\end{definition}

Para o caminho \omterm{def:b3:pericyclic:rotation-modes}{disrotatório} proibido, o estado fundamental do butadieno ($\psi_1^2\psi_2^2$)
se correlaciona com um estado duplamente excitado do ciclobuteno ($\sigma^2\pi^{*2}$), mas o primeiro estado
excitado ($\psi_1^2\psi_2\psi_3$) se correlaciona com o primeiro estado excitado do ciclobuteno ($\sigma^2\pi\pi^*$): a partir do
estado excitado, o caminho \omterm{def:b3:pericyclic:rotation-modes}{disrotatório} é descendente. É a inversão fotoquímica
da \cref{prop:b3:pericyclic:photo-electrocyclic}, vista nos estados.

\section{Cicloadições}

\begin{theorem}[Cicloadições]\label{thm:b3:pericyclic:cycloaddition}
Uma \omterm{def:b3:pericyclic:families}{cicloadição} $[m + n]$ térmica na qual os dois componentes reagem de modo \omterm{def:b3:pericyclic:faciality}{suprafacial} é
permitida quando $m + n = 4q + 2$ elétrons, e proibida quando $m + n = 4q$; sob luz, a
regra se inverte.
\end{theorem}

\begin{proof}
As duas novas ligações se formam entre a HOMO de um componente e a LUMO do
outro, que se defrontam por seus lobos terminais; com os dois reagindo
de modo \omterm{def:b3:pericyclic:faciality}{suprafacial}, as duas ligações só são ligantes se os coeficientes terminais da
HOMO e da LUMO têm o mesmo sinal relativo. Para um componente de $N$
elétrons em $N$ átomos, a HOMO tem $k = N/2$ e a LUMO $k = N/2 + 1$; coeficientes terminais de
mesmo sinal para $k$ ímpar, opostos para $k$ par. No caso [4+2], a HOMO do dieno ($k = 2$)
e a LUMO do eteno ($k = 2$) ambas têm sinais terminais opostos: elas se correspondem. No
caso [2+2], a HOMO de um eteno ($k = 1$, mesmos sinais) encontra a LUMO do outro
($k = 2$, opostos): uma extremidade ligante, a outra antiligante. Em geral, a HOMO do componente $m$
e a LUMO do componente $n$ se correspondem quando $m/2$ e $n/2 + 1$ têm a mesma
paridade, isto é, quando $(m + n)/2$ é ímpar: $m + n = 4q + 2$. Sob luz, a LUMO monoocupada do
componente excitado faz o papel de sua HOMO, e a condição de paridade se
inverte.
\end{proof}

\begin{omfigure}
\begin{tikzpicture}[font=\scriptsize]
  % [4+2]: diene HOMO above, ethene LUMO (pi*) below, facing lobes match at both ends
  \foreach \x/\c in {0/0.60, 1/0.37, 2/-0.37, 3/-0.60} {
    \pic at (\x,1.6) {omp={1.4*\c}{90}}; \fill (\x,1.6) circle (1.1pt);
  }
  \draw[black!60] (0,1.6) -- (3,1.6);
  \foreach \x/\c in {0/-0.71, 3/0.71} {
    \pic at (\x,-0.4) {omp={1.4*\c}{90}}; \fill (\x,-0.4) circle (1.1pt);
  }
  \draw[black!60] (0,-0.4) -- (3,-0.4);
  \draw[omMeth, densely dashed, thick] (-0.25,0.75) -- (-0.25,0.35) (3.25,0.75) -- (3.25,0.35);
  \node[omMeth] at (-0.25,0.55) [left] {ligante};
  \node[omMeth] at (3.25,0.55) [right] {ligante};
  \node[font=\small] at (1.5,2.95) {[4+2]: HOMO do dieno};
  \node[font=\small] at (1.5,-1.75) {LUMO do eteno: permitida};
  \begin{scope}[xshift=7.2cm]
  \foreach \x/\c in {0/0.71, 3/0.71} {
    \pic at (\x,1.6) {omp={1.4*\c}{90}}; \fill (\x,1.6) circle (1.1pt);
  }
  \draw[black!60] (0,1.6) -- (3,1.6);
  \foreach \x/\c in {0/-0.71, 3/0.71} {
    \pic at (\x,-0.4) {omp={1.4*\c}{90}}; \fill (\x,-0.4) circle (1.1pt);
  }
  \draw[black!60] (0,-0.4) -- (3,-0.4);
  \draw[omMeth, densely dashed, thick] (-0.25,0.75) -- (-0.25,0.35);
  \draw[omThm, densely dashed, thick] (3.25,0.75) -- (3.25,0.35);
  \node[omMeth] at (-0.25,0.55) [left] {ligante};
  \node[omThm] at (3.25,0.55) [right] {antiligante};
  \node[font=\small] at (1.5,2.95) {[2+2]: HOMO do eteno};
  \node[font=\small] at (1.5,-1.75) {LUMO do eteno: proibida};
  \end{scope}
\end{tikzpicture}

{\small Orbitais de fronteira frente a frente em \omterm{def:b3:pericyclic:families}{cicloadições} \omterm{def:b3:pericyclic:faciality}{suprafaciais} (os
lobos inferiores do componente de cima encontram os lobos superiores do de baixo; tamanho
dos lobos segundo os coeficientes de Hückel). Na reação [4+2], as duas novas ligações são
ligantes; na reação [2+2], uma extremidade é ligante e a outra antiligante.}
\end{omfigure}

Termicamente, dois alquenos não dão um ciclobutano em uma única etapa concertada; sob
luz, dão (\cref{ch:b3:photochemistry}). As cetenas são a exceção que
confirma a regra: seu orbital $\pi^*$ C=O, perpendicular ao C=C, permite que elas
reajam com um alqueno em uma geometria [$\pi2_s$ + $\pi2_a$], um componente \omterm{def:b3:pericyclic:faciality}{suprafacial} e a
cetena \omterm{def:b3:pericyclic:faciality}{antarafacial}, e elas formam ciclobutanonas termicamente.

\begin{definition}[Cicloadições 1,3-dipolares]\label{def:b3:pericyclic:dipolar}
Um \emph{1,3-dipolo}\index{1,3-dipolo} é um sistema $\pi$ de três átomos com quatro
elétrons, escrito com cargas formais nas extremidades ou no meio, como uma
azida $\ce{R-N=N+=N-}$, um óxido de nitrila ou uma molécula de ozônio. Uma \emph{cicloadição
1,3-dipolar}\index{cicloadicao 1,3-dipolar@cicloadição 1,3-dipolar} é sua \omterm{def:b3:pericyclic:families}{cicloadição} [3+2], uma
reação de seis elétrons, com um alqueno ou um alquino, que dá um anel de cinco membros.
\end{definition}

A reação, catalisada por cobre, de uma azida com um alquino terminal, que dá um único
regioisômero de um 1,2,3-triazol, é a reação de junção confiável da química
click. A ozonólise começa com uma \omterm{def:b3:pericyclic:dipolar}{cicloadição 1,3-dipolar} do ozônio ao
alqueno.

\section{Rearranjos sigmatrópicos e a regra geral}

\begin{proposition}[Deslocamentos de hidrogênio]\label{prop:b3:pericyclic:sigmatropic-h}
Um deslocamento $[1, j]$ térmico de hidrogênio ao longo de um polieno é permitido de modo \omterm{def:b3:pericyclic:faciality}{suprafacial} quando o
estado de transição contém $4q + 2$ elétrons ($[1,5]$) e precisa ser \omterm{def:b3:pericyclic:faciality}{antarafacial} quando contém $4q$
($[1,3]$, $[1,7]$).
\end{proposition}

\begin{proof}
Trate o estado de transição como um átomo de hidrogênio (1s) que se move entre as extremidades de um
radical de $j$ átomos com $j$ elétrons; o hidrogênio faz ponte entre as duas extremidades da
SOMO, $k = (j + 1)/2$. Para $j = 5$, $k = 3$ é ímpar, as extremidades têm lobos de mesmo sinal em uma face,
e o orbital 1s pode se ligar a ambas nessa face: \omterm{def:b3:pericyclic:faciality}{suprafacial}. Para $j = 3$ e $j = 7$,
$k = 2$ e $4$ são pares, os lobos de mesmo sinal ficam em faces opostas nas duas extremidades, e
o hidrogênio precisa passar de uma face para a outra: \omterm{def:b3:pericyclic:faciality}{antarafacial}. Um deslocamento $[1,3]$
não consegue alcançar através de uma cadeia tão curta; um deslocamento $[1,7]$ consegue, ao longo de um polieno helicoidal.
\end{proof}

\begin{definition}[Rearranjos de Cope e de Claisen]\label{def:b3:pericyclic:cope-claisen}
O \emph{rearranjo de Cope}\index{rearranjo de Cope} é o \omterm{def:b3:pericyclic:families}{rearranjo sigmatrópico} [3,3]
de um hexa-1,5-dieno; o \emph{rearranjo de
Claisen}\index{rearranjo de Claisen}, o de um alil vinil éter, ou de
um alil aril éter, em um composto carbonílico $\gamma,\delta$-insaturado ou em um \textit{orto}-alilfenol.
\end{definition}

\begin{omfigure}
\begin{tikzpicture}[font=\scriptsize, scale=1.1]
  \foreach \px/\three/\title in {0/C/Cope, 5.0/O/Claisen} {
    \begin{scope}[xshift=\px cm]
      % chair: atoms at 60-degree steps, alternately 0.25 above and below the
      % mean plane, seen from slightly above (y' = 0.35 y + z), scaled
      \begin{scope}[shift={(1.3,0.2)}, scale=1.15]
      \coordinate (a1) at (-1.25,-0.25); \coordinate (a2) at (-0.625,-0.129); \coordinate (a3) at (0.625,-0.629);
      \coordinate (a4) at (1.25,0.25); \coordinate (a5) at (0.625,0.129); \coordinate (a6) at (-0.625,0.629);
      \end{scope}
      \draw[thick] (a1) -- (a2) -- (a3);
      \draw[thick] (a4) -- (a5) -- (a6);
      \draw[omThm, thick, densely dashed] (a3) -- (a4);
      \draw[omMeth, thick, densely dashed] (a6) -- (a1);
      \foreach \i/\l in {1/1, 2/2, 4/4, 5/5, 6/6} { \fill (a\i) circle (1.3pt); }
      \node[fill=white, inner sep=0.6pt] at (a3) {\three};
      \node[anchor=east] at (a1) {1};
      \node[anchor=north] at (a2) {2};
      \node[anchor=north west] at (a3) {3};
      \node[anchor=west] at (a4) {4};
      \node[anchor=south] at (a5) {5};
      \node[anchor=south] at (a6) {6};
      \node[font=\small\bfseries] at (1.3,1.5) {\title};
    \end{scope}
  }
\end{tikzpicture}

{\small Estados de transição em cadeira de \omterm{def:b3:pericyclic:families}{rearranjos sigmatrópicos} [3,3]: dois fragmentos
alila se defrontam; a ligação $\sigma$ 3--4 se rompe (vermelho) enquanto a ligação 1--6 se forma
(verde). No \omterm{def:b3:pericyclic:cope-claisen}{rearranjo de Claisen}, o átomo 3 é o oxigênio do éter, e o
produto é um composto carbonílico. A cadeira fixa a configuração relativa dos
novos estereocentros.}
\end{omfigure}

O deslocamento [3,3] é um processo de seis elétrons, inteiramente \omterm{def:b3:pericyclic:faciality}{suprafacial}, termicamente permitido;
ele passa por um estado de transição do tipo cadeira, como um ciclo-hexano, e
os substituintes preferem suas posições equatoriais: os dienos (\textit{E},\textit{E}) e (\textit{Z},\textit{Z}) dão
configurações relativas opostas, de modo previsível.

\begin{theorem}[Regras de Woodward--Hoffmann]\label{thm:b3:pericyclic:woodward-hoffmann}
Uma \omterm{def:b3:pericyclic:pericyclic}{reação pericíclica} térmica é permitida quando o número total de componentes $(4q + 2)_s$ e
$(4r)_a$ é ímpar; uma fotoquímica, quando é par.
\end{theorem}

\begin{proof}[Demonstração parcial]
Para cada família tratada acima, a contagem concorda com a análise
dos orbitais de fronteira: uma \omterm{def:b3:pericyclic:families}{cicloadição} [4+2] térmica tem componentes $\pi4_s$ e $\pi2_s$, um componente
$(4q + 2)_s$, ímpar: permitida; a [2+2] tem $\pi2_s + \pi2_s$, dois, par: proibida, ao passo que
$\pi2_s + \pi2_a$ tem um: permitida; a abertura \omterm{def:b3:pericyclic:rotation-modes}{conrotatória} do ciclobuteno é $\sigma2_s + \pi2_a$, um:
permitida; o deslocamento $[1,5]$-H \omterm{def:b3:pericyclic:faciality}{suprafacial} é $\sigma2_s + \pi4_s$, onde só $\sigma2_s$ conta: um,
permitido. A demonstração geral, válida para todo arranjo de componentes, pertence
a cursos mais avançados.
\end{proof}

\begin{definition}[Estados de transição aromáticos e de Möbius]\label{def:b3:pericyclic:mobius}
Um \emph{estado de transição aromático}\index{estado de transicao aromatico@estado de transição aromático} é um arranjo
cíclico de orbitais em um estado de transição pericíclico que, como um anel de Hückel, é
estabilizado quando contém $4q + 2$ elétrons. Um \emph{estado de transição de
Möbius}\index{estado de transicao de Mobius@estado de transição de Möbius} tem um número ímpar de mudanças de sinal entre
orbitais vizinhos ao redor do anel (uma torção \omterm{def:b3:pericyclic:faciality}{antarafacial}); ele é
estabilizado com $4q$ elétrons.
\end{definition}

\begin{proposition}[Estados de transição aromáticos]\label{prop:b3:pericyclic:aromatic-ts}
Uma \omterm{def:b3:pericyclic:pericyclic}{reação pericíclica} térmica é permitida quando seu estado de transição é aromático:
topologia de Hückel com $4q + 2$ elétrons, ou topologia de Möbius com $4q$ elétrons.
\end{proposition}

\begin{proof}[Argumento]
O arranjo cíclico de orbitais no estado de transição se comporta como uma molécula
conjugada cíclica. Um anel de Hückel tem seu nível mais baixo não degenerado e pares
acima dele, preenchidos com $4q + 2$ elétrons (volume do 2.º ano); um anel com uma inversão
de sinal tem todos os seus níveis em pares, preenchidos com $4q$. Uma camada preenchida significa
estabilização, um \omterm{def:b3:pericyclic:mobius}{estado de transição aromático} e uma barreira baixa. A contagem das mudanças
de sinal é a contagem dos componentes \omterm{def:b3:pericyclic:faciality}{antarafaciais}, de modo que essa regra é equivalente à
regra de Woodward--Hoffmann.
\end{proof}

\begin{method}[Desenhar e contar os componentes]\label{met:b3:pericyclic:components}
\begin{enumerate}
  \item Marque as ligações formadas e rompidas; cada sistema $\pi$ ou ligação $\sigma$ que muda é um
        componente, rotulado com seu tipo e sua contagem de elétrons ($\pi4$, $\pi2$, $\sigma2$, $\omega0$ para
        um orbital vazio).
  \item Escolha s ou a para cada componente a partir da geometria (mesma face, ou
        faces opostas).
  \item Conte os componentes $(4q + 2)_s$ e $(4r)_a$: ímpar, permitida termicamente; par, permitida
        fotoquimicamente.
  \item Confira com a contagem de elétrons ao redor do anel: Hückel com $4q + 2$,
        Möbius com $4q$.
\end{enumerate}
\end{method}

\begin{omfigure}
\begin{tikzpicture}[font=\scriptsize, >=Stealth]
  \foreach \px/\mode in {0/1, 4.6/2, 9.2/3} {
    \begin{scope}[xshift=\px cm]
      \foreach \n/\a in {10/210, 5/150, 6/90, 7/30, 8/-30, 9/-90} { \coordinate (v\n) at (\a:0.85); }
      \ifnum\mode=1
        \draw[thick] (v10) -- (v5) (v6) -- (v7) (v8) -- (v9);
        \draw[thick, double, double distance=1.6pt] (v5) -- (v6) (v7) -- (v8);
        \draw[thick, omThm] (v9) -- (v10);
        \draw[thick] (v10) -- ++(240:0.6) node[anchor=north] {$\ce{CH3}$};
        \node[font=\small\bfseries] at (0,1.55) {7-desidrocolesterol};
      \fi
      \ifnum\mode=2
        \draw[thick, double, double distance=1.6pt] (v10) -- (v5) (v6) -- (v7) (v8) -- (v9);
        \draw[thick] (v5) -- (v6) (v7) -- (v8);
        \draw[thick] (v10) -- ++(240:0.6) node[anchor=north] (me) {$\ce{CH3}$};
        \draw[->, omThm, thick] ($(v10)+(240:0.75)+(0.35,0)$) to[bend right=40] ($(v9)+(-0.12,-0.12)$);
        \node[font=\small\bfseries] at (0,1.55) {pré-vitamina D$_3$};
      \fi
      \ifnum\mode=3
        \draw[thick] (v10) -- (v5) (v6) -- (v7) (v8) -- (v9);
        \draw[thick, double, double distance=1.6pt] (v5) -- (v6) (v7) -- (v8);
        \draw[thick, double, double distance=1.6pt] (v10) -- ++(240:0.6) node[anchor=north] {$\ce{CH2}$};
        \node[anchor=north west, inner sep=1pt] at (v9) {$\ce{CH2}$};
        \node[font=\small\bfseries] at (0,1.55) {vitamina D$_3$};
      \fi
      \foreach \n/\a in {10/210, 5/150, 6/90, 7/30, 8/-30, 9/-110} {
        \node[font=\tiny, black!60] at (\a:1.12) {\n};
      }
    \end{scope}
  }
  \draw[->, thick] (1.45,0) -- (3.0,0) node[midway, above, align=center] {$h\nu$\\ 6 e$^-$, conrotatório};
  \draw[->, thick] (6.05,0) -- (7.6,0) node[midway, above, align=center] {calor, $[1,7]$-H\\ antarafacial};
\end{tikzpicture}

{\small As duas etapas pericíclicas da síntese da vitamina D, desenhadas sobre os átomos
que reagem (numeração esteroide em cinza; o restante da molécula, anéis A, C e D,
omitido). A luz abre o anel B rompendo a ligação 9--10 (vermelho), uma \omterm{def:b3:pericyclic:families}{reação eletrocíclica}
fotoquímica de seis elétrons, \omterm{def:b3:pericyclic:rotation-modes}{conrotatória}. Na pré-vitamina D$_3$, um hidrogênio do
grupo metila C19 passa então para C9 (seta vermelha) através do trieno
helicoidal de sete átomos, de modo \omterm{def:b3:pericyclic:faciality}{antarafacial}.}
\end{omfigure}

\begin{inthelab}[Um rearranjo de Claisen]
Um alil aril éter é aquecido a refluxo, sob nitrogênio, em um solvente de alto ponto de
ebulição, como o 1,2-diclorobenzeno, ou sem solvente, a cerca de \qty{200}{\celsius}, e
a reação é acompanhada por cromatografia em camada delgada: o \textit{orto}-alilfenol formado
é mais polar que o éter. O fenol é então extraído em hidróxido de sódio
aquoso, separado das impurezas neutras e recuperado por acidificação.
\end{inthelab}

\begin{safety}{\ghs[0.62cm]{GHS02}\ghs[0.62cm]{GHS04}\ghs[0.62cm]{GHS08}}
% ledger: ghs4:butadiene
Buta-1,3-dieno, um dieno comum: gás extremamente inflamável sob pressão, pode provocar
defeitos genéticos e câncer. É manipulado somente em sistemas fechados; no ensino,
usa-se em seu lugar, em aparelhagem fechada, o sulfoleno sólido, que o libera em uma \omterm{def:b3:pericyclic:families}{reação queletrópica}
quando aquecido.
\end{safety}

\begin{history}[A simetria orbital]
Em 1965, Robert Woodward, que encontrara a intrigante estereoquímica de
etapas eletrocíclicas ao sintetizar a vitamina B$_{12}$, e Roald Hoffmann publicaram
as regras de conservação da simetria orbital. Kenichi Fukui introduzira os
orbitais de fronteira em 1952; Fukui e Hoffmann dividiram o Prêmio Nobel de
Química de 1981. Woodward, que recebera o prêmio de 1965 por suas sínteses, morreu em
1979.
\end{history}

\section{Exercícios}

\begin{exercise}[$\star$]\label{exo:b3:pericyclic:1}
Classifique: a reação de Diels--Alder; a abertura do anel do ciclobuteno; o \omterm{def:b3:pericyclic:cope-claisen}{rearranjo de
Cope}; a perda de $\ce{SO2}$ pelo sulfoleno; a adição do ozônio a um alqueno;
um deslocamento $[1,5]$-H no ciclopentadieno.
\end{exercise}

\begin{exercise}[$\star$]\label{exo:b3:pericyclic:2}
Preveja o modo, \omterm{def:b3:pericyclic:rotation-modes}{conrotatório} ou \omterm{def:b3:pericyclic:rotation-modes}{disrotatório}, do fechamento de anel térmico e do
fotoquímico do hexa-1,3,5-trieno e do buta-1,3-dieno.
\end{exercise}

\begin{exercise}[$\star$]\label{exo:b3:pericyclic:3}
Que produto o \textit{trans}-3,4-dimetilciclobuteno dá por aquecimento?
\end{exercise}

\begin{exercise}[$\star$]\label{exo:b3:pericyclic:4}
Por que o eteno não se dimeriza a ciclobutano por aquecimento, mas o faz sob luz
(com um sensibilizador ou por excitação direta)?
\end{exercise}

\begin{exercise}[$\star\star$]\label{exo:b3:pericyclic:5}
Construa o \omterm{def:b3:pericyclic:orbital-correlation}{diagrama de correlação de orbitais} de hexatrieno $\rightleftharpoons$ ciclo-hexa-1,3-dieno para o
caminho \omterm{def:b3:pericyclic:rotation-modes}{disrotatório}, e conclua.
\end{exercise}

\begin{exercise}[$\star\star$]\label{exo:b3:pericyclic:6}
Sob luz, o (\textit{E},\textit{Z},\textit{E})-octa-2,4,6-trieno se fecha em qual dimetilciclo-hexadieno?
\end{exercise}

\begin{exercise}[$\star\star$]\label{exo:b3:pericyclic:7}
O 5-metilciclopenta-1,3-dieno se rearranja à temperatura ambiente em 1- e
2-metilciclopentadienos. Explique por um deslocamento $[1,5]$-H e diga por que ele é \omterm{def:b3:pericyclic:faciality}{suprafacial}.
\end{exercise}

\begin{exercise}[$\star\star$]\label{exo:b3:pericyclic:8}
Conte os componentes da \omterm{def:b3:pericyclic:families}{cicloadição} de uma cetena a um alqueno e mostre
que ela é termicamente permitida.
\end{exercise}

\begin{exercise}[$\star\star$]\label{exo:b3:pericyclic:9}
O alil vinil éter, $\ce{CH2=CH-O-CH2-CH=CH2}$, sofre o \omterm{def:b3:pericyclic:cope-claisen}{rearranjo de Claisen} por
aquecimento. Desenhe seu estado de transição em cadeira e dê o produto.
\end{exercise}

\begin{exercise}[$\star\star\star$]\label{exo:b3:pericyclic:10}
Um estado de transição tem oito elétrons em um arranjo cíclico com um componente
\omterm{def:b3:pericyclic:faciality}{antarafacial}. Ele é aromático? A reação é termicamente permitida?
\end{exercise}

\begin{exercise}[$\star\star\star$]\label{exo:b3:pericyclic:11}
Na \omterm{def:b3:pericyclic:dipolar}{cicloadição 1,3-dipolar} de uma azida com um alquino terminal, o maior
coeficiente da HOMO da azida fica no nitrogênio terminal N3 e o maior
coeficiente da LUMO do alquino, em seu carbono terminal (dados do exercício).
Preveja o regioisômero favorecido sem catalisador emparelhando os maiores
coeficientes, e comente.
\end{exercise}

\begin{exercise}[$\star\star\star$]\label{exo:b3:pericyclic:12}
Mostre, usando a fórmula de Hückel, que o deslocamento $[1,7]$-H precisa ser \omterm{def:b3:pericyclic:faciality}{antarafacial},
e o deslocamento $[1,5]$-H, \omterm{def:b3:pericyclic:faciality}{suprafacial}.
\end{exercise}

% keep the heading with its problem box, which fills a page by itself
\clearpage
\enlargethispage{3\baselineskip}
\section{Problema: Sol, pele e vitamina D}

\begin{problem}[{Problema de fim de semana --- sol, pele e vitamina D: a
abertura de anel fotoquímica do 7-desidrocolesterol, o deslocamento térmico de hidrogênio,
os subprodutos e quanto tempo o organismo leva para terminar o trabalho}]
\label{pb:b3:pericyclic:1}
Dados do problema: a \qty{37}{\celsius}, a conversão da pré-vitamina D$_3$ em vitamina D$_3$ é
de primeira ordem com $k = \qty{2.3e-2}{h^{-1}}$ (sua reação inversa é desprezada); sua energia de ativação é
\qty{85}{kJ/mol}.

\medskip\noindent\textbf{Parte I --- A abertura do anel.}
\begin{enumerate}
  \item Que ligação do 7-desidrocolesterol se rompe, e que sistema $\pi$ se forma?
  \item Quantos elétrons participam, e qual é a família da reação?
  \item Por que ela precisa de luz?
  \item A abertura fotoquímica é \omterm{def:b3:pericyclic:rotation-modes}{conrotatória} ou \omterm{def:b3:pericyclic:rotation-modes}{disrotatória}? Justifique com os
        orbitais.
  \item Por que a nova ligação dupla central da pré-vitamina D$_3$ precisa ser \textit{Z}?
  \item Por que a mesma abertura não ocorre termicamente, no escuro, no corpo?
\end{enumerate}

\medskip\noindent\textbf{Parte II --- O deslocamento de hidrogênio.}
\begin{enumerate}[resume]
  \item Nomeie os átomos entre os quais o hidrogênio se move, e o tipo de deslocamento.
  \item Quantos elétrons participam do estado de transição?
  \item Conte os componentes: a reação térmica é permitida
        de modo \omterm{def:b3:pericyclic:faciality}{suprafacial} ou \omterm{def:b3:pericyclic:faciality}{antarafacial}.
  \item Por que um deslocamento \omterm{def:b3:pericyclic:faciality}{antarafacial} é possível aqui, e um deslocamento $[1,3]$
        não?
  \item O estado de transição é de Hückel ou de Möbius, e é aromático?
  \item Por que o deslocamento não precisa de luz?
\end{enumerate}

\medskip\noindent\textbf{Parte III --- Subprodutos.}
\begin{enumerate}[resume]
  \item A pré-vitamina D$_3$ pode voltar a se fechar sob luz. Em que modo, e por que isso pode dar um
        estereoisômero de anel do 7-desidrocolesterol (o lumisterol)?
  \item A luz também pode transformar a ligação dupla central \textit{Z} em \textit{E} (taquisterol). Por que
        o taquisterol não pode sofrer o deslocamento $[1,7]$-H?
  \item Por que uma exposição prolongada ao sol não produz cada vez mais vitamina D?
  \item Por que a abertura do anel é rápida, mas o deslocamento de hidrogênio é lento?
\end{enumerate}

\medskip\noindent\textbf{Parte IV --- Cinética à temperatura do corpo.}
\begin{enumerate}[resume]
  \item Calcule a meia-vida da pré-vitamina D$_3$ a \qty{37}{\celsius}.
  \item Calcule a fração convertida em um dia.
  \item Calcule o tempo para converter 90\,\%.
  \item Calcule $k$ a \qty{20}{\celsius}.
  \item Calcule o tempo para converter 90\,\% a \qty{20}{\celsius}.
  \item Por que é útil que a pele seja quente?
  \item Por que os suplementos de vitamina D são feitos irradiando e aquecendo um
        esterol, e não por síntese clássica?
  \item Enuncie o resultado: o tempo para converter 90\,\% da pré-vitamina D$_3$ no corpo.
\end{enumerate}
\end{problem}
